% !TeX root = ../paper.tex
\section{SmartVLA-Pi and SmartVLA-OFT: Architecture, Training, and Calling Protocol}
\label{appendix:model-details}

This appendix describes the two policy instantiations used in the
experiments. SmartVLA changes the supervision and training procedure,
while retaining the native action models of the two base VLA families.
Both variants consume the same type of caller-provided instruction,
third-person and wrist RGB observations, and robot state, and return a
finite local action block. The regression and flow-matching losses in
Equations~(\ref{eq:oft-loss}) and~(\ref{eq:pi-flow-loss}) supply the native
losses used in the shared training formulation of Equation~(\ref{eq:action-objective}).

\subsection{Common Instruction-Conditioned Interface}

At time $t$, each policy maps the current multimodal observation and
command to a local action distribution,
\begin{equation}
\label{eq:policy-interface}
\mathbf{a}_{t:t+H-1} \sim \pi_\theta(\cdot \mid
  \mathbf{o}_t,\mathbf{s}_t,c_t),
\end{equation}
where $c_t$ can be a complete task instruction, a phase-level
subgoal, an atomic motion, a spatially grounded command, or a corrective
instruction. The visual markers used by LIBERO-SMART are rendered in the
RGB inputs and are not supplied through privileged channels. The interface in
Equation~(\ref{eq:policy-interface}) is used for end-to-end task execution and for repeated calls
from an upper-level controller.

\subsection{SmartVLA-OFT}
\label{app:smartvla-oft}

SmartVLA-OFT uses Qwen3-VL-4B-Instruct~\citep{bai2025qwen3vl} as its vision-language backbone
and an OFT-style residual MLP for continuous action regression. The
backbone combines a 24-layer vision transformer with a 36-layer language
transformer of hidden width 2,560. Multi-view image tokens, the current
instruction, and a textual encoding of normalized robot state form the
conditioning sequence. We append $H$ action-query tokens, one for each
predicted step. The final-layer hidden states at these positions are
extracted together in a single forward pass:
\begin{equation}
\label{eq:oft-features}
\mathbf{Z}_t=\operatorname{Gather}_{\mathrm{act}}
\!\left[F_\psi(\mathbf{o}_t,c_t,
\operatorname{Tok}_s(\mathbf{s}_t),\mathbf{q}_{1:H})\right]
\in\mathbb{R}^{H\times2560}.
\end{equation}
In Equation~(\ref{eq:oft-features}), $\operatorname{Tok}_s$ denotes state discretization and textual
encoding, and $\mathbf{q}_{1:H}$ are fixed action placeholders rather
than previously generated actions. A shared MLP maps each query state
to a normalized continuous action. It uses a 5,120-dimensional hidden
representation, two residual blocks with LayerNorm and ReLU, and a
32-dimensional output projection:
\begin{equation}
\label{eq:oft-action-decoding}
\widehat{\mathbf{A}}_t=g_\phi(\mathbf{Z}_t),\qquad
\mathbf{A}^{\mathrm{exec}}_t=
\mathcal{N}^{-1}\!\left(\widehat{\mathbf{A}}_t[:,\mathcal{I}]\right).
\end{equation}
In Equation~(\ref{eq:oft-action-decoding}), the index set $\mathcal{I}$ selects the active robot
action dimensions,
and $\mathcal{N}^{-1}$ applies the corresponding action denormalization.
Training uses an $L_1$ loss on these active dimensions:
\begin{equation}
\label{eq:oft-loss}
\mathcal{L}_{\mathrm{OFT}}=
\frac{1}{H|\mathcal{I}|}\sum_{k=0}^{H-1}\sum_{j\in\mathcal{I}}
\left|\widehat{A}_{t,k,j}-A_{t,k,j}\right|.
\end{equation}
The LIBERO instantiation uses two $256\times256$ RGB views, an
eight-dimensional state encoded on a 0--1000 scale, $H=8$, and seven
active action dimensions. The whole chunk is predicted at once, without
autoregressive action-token decoding or iterative flow sampling.
Action-aligned segments and grade-conditioned rewrites provide the
semantic and spatial supervision for this otherwise unchanged
continuous-regression interface.

\subsection{SmartVLA-Pi}
\label{app:smartvla-pi}

SmartVLA-Pi uses the $\pi_{0.5}$ architecture implemented in OpenPI:
a PaliGemma vision-language backbone comprising a SigLIP-So400m/14
visual encoder and a Gemma-2B language model, paired with a separate
approximately 300M-parameter action expert. The expert has 18 transformer
layers and hidden width 1,024. Images, language, and discretized robot
state form the multimodal prefix. In this $\pi_{0.5}$ configuration,
state is encoded with the language tokens rather than injected as a
separate continuous state token into the action suffix. Noisy actions
are projected into expert tokens, while a sinusoidal flow-time embedding
and an MLP condition the expert through adaptive RMS normalization.
Expert tokens attend to the multimodal prefix and to one another.

Let $\mathbf{A}_t\in\mathbb{R}^{H\times32}$ denote the normalized
demonstration action chunk and $\epsilon\sim\mathcal{N}(0,I)$ Gaussian
noise. Using the convention that $\tau=1$ is noise and $\tau=0$ is the
action target, training constructs
\begin{equation}
\label{eq:pi-flow-interpolation}
\mathbf{X}^{\tau}_t=(1-\tau)\mathbf{A}_t+\tau\epsilon,
\qquad \mathbf{u}_t=\epsilon-\mathbf{A}_t.
\end{equation}
Let $\mathcal{I}$ denote the active action dimensions of the current
embodiment and $\mathbf{M}$ the binary mask retaining those dimensions
at every predicted step. Using the interpolation and velocity target in
Equation~(\ref{eq:pi-flow-interpolation}), the conditional velocity field $v_\theta$
is learned through a flow-matching loss evaluated only on active dimensions,
\begin{equation}
\label{eq:pi-flow-loss}
\mathcal{L}_{\mathrm{Pi}}=
\mathbb{E}_{\mathbf{A}_t,\epsilon,\tau}\!
\left[\frac{1}{H|\mathcal{I}|}\left\|
\mathbf{M}\odot\left(
v_\theta(\mathbf{X}^{\tau}_t,\tau;\mathbf{P}_t)
-(\epsilon-\mathbf{A}_t)\right)\right\|_F^2\right],
\end{equation}
where $\mathbf{P}_t$ denotes the encoded multimodal prefix. Padded
dimensions are excluded from the loss in warmup, pretraining, and adaptation alike.
At inference, the prefix is encoded once and its attention keys and
values are cached. Starting from
$\mathbf{X}^{(0)}_t=\mathbf{M}\odot\epsilon$, the action expert performs
$N$ Euler steps from noise to actions:
\begin{equation}
\label{eq:pi-flow-inference}
\begin{aligned}
\tau_n&=1-\frac{n}{N},\qquad n=0,\ldots,N-1,\\
\mathbf{X}^{(n+1)}_t&=\mathbf{M}\odot\left[
\mathbf{X}^{(n)}_t-\frac{1}{N}v_\theta(\mathbf{X}^{(n)}_t,\tau_n;\mathbf{P}_t)\right],\\
\mathbf{A}^{\mathrm{exec}}_t&=
\mathcal{N}^{-1}\!\left(\mathbf{X}^{(N)}_t[:,\mathcal{I}]\right).
\end{aligned}
\end{equation}
The LIBERO instantiation uses $224\times224$ model-input images,
$H=10$, $N=10$, and seven active dimensions in the 32-dimensional
action container; inactive dimensions remain zero during sampling.
In Equation~(\ref{eq:pi-flow-inference}), $N$ counts internal integration steps, whereas $H$ counts
predicted environment actions. Prefix caching avoids rerunning the
vision-language backbone at every integration step. SmartVLA changes
the instruction and segment-level action supervision while retaining this
flow-matching action generator.

\subsection{Three-Stage Optimization}\label{app:training-hparams}

Both variants follow the same three-stage schedule: warmup,
pretraining, and adaptation. Warmup uses object-free atomic-motion commands and local action
segments. With images omitted and robot state retained, the entire VLM
backbone is frozen while the action head learns the correspondence between
a short command and its local control. Pretraining restores the multimodal observations and jointly
trains the policy on original instructions and grade-permitted semantic
and spatial rewrites. Adaptation then fine-tunes the
complete policy independently for each target domain, mixing original
task commands with the applicable steering supervision. Original-task
examples preserve the continuous episode trajectory, whereas local
examples use the corresponding action-aligned segment. This separates
complete-task imitation from local command following while keeping both
forms accessible through one instruction interface.

We collectively refer to the checkpoints obtained after warmup
followed by pretraining as \textbf{SmartVLA-Base}, with two instantiations:
\textbf{SmartVLA-Pi-Base} and \textbf{SmartVLA-OFT-Base}. The warmup checkpoint
is an intermediate initialization for pretraining. For each architecture, its
Base checkpoint initializes independent adaptation runs for LIBERO, RoboTwin
Full, RoboTwin C2R, and Piper. Each run produces a separate adapted expert;
adaptation does not proceed sequentially from one target-domain expert to
another. Unless otherwise specified, SmartVLA results use the corresponding
adapted experts, reported as SmartVLA-Pi and SmartVLA-OFT. The Base suffix
identifies the pretrained checkpoint before target-domain adaptation.

SmartVLA-Pi and SmartVLA-OFT use the same listed training settings and update
counts within each stage and target-domain protocol. Shared warmup and pretraining
settings are reported once in Table~\ref{tab:pretraining-hparams}, with
stage-specific learning rates indicated explicitly. Adaptation settings are
listed by target domain in Table~\ref{tab:domain-finetuning-hparams}.
Training steps count optimizer updates, and the effective global batch size
includes data parallelism and gradient accumulation.

% !TeX root = ../paper.tex
\begin{center}
\begin{minipage}{\linewidth}
\centering
\normalsize
\setlength{\tabcolsep}{3pt}
\renewcommand{\arraystretch}{1.10}
\setlength{\belowcaptionskip}{\baselineskip}
\captionof{table}{Shared warmup and pretraining hyperparameters for SmartVLA-Pi and SmartVLA-OFT. Stage-specific values are indicated explicitly.}\label{tab:pretraining-hparams}
\begin{tabular}{@{}>{\raggedright\arraybackslash}p{0.460\linewidth}>{\raggedright\arraybackslash}p{0.490\linewidth}@{}}
\toprule
Parameter & Shared configuration \\
\midrule
Trainable modules during warmup & Action head; VLM backbone frozen \\
Trainable modules during pretraining & All components; no freezing \\
Original/local instruction sampling (pretraining and adaptation) & 1:1 effective probability; not applicable to warmup \\
Hardware & $16\times$ NVIDIA A100 \\
Effective global batch size & 256 \\
Optimizer & Fused AdamW \\
Peak learning rate (warmup / pretraining) & $5\times10^{-5}$ / $10^{-5}$ \\
Adam coefficients $(\beta_1,\beta_2)$ & $(0.9,0.95)$ \\
Optimizer epsilon & $10^{-8}$ \\
Weight decay & 0.1 \\
Learning-rate schedule & \texttt{cosine\_with\_min\_lr} \\
Learning-rate warmup ratio & 0.01 \\
Gradient-clipping threshold & 1.0 \\
TF32 & Enabled \\
Parallelism & DeepSpeed ZeRO-2 \\
State normalization (\texttt{state\_mode}) & \texttt{QUANTILE} (Q01--Q99) \\
Training random seed & 42 \\
Optimizer updates (warmup / pretraining) & 10,000 / 240,000 \\
\bottomrule
\end{tabular}
\end{minipage}
\end{center}
%

% !TeX root = ../paper.tex
\begin{center}
\begin{minipage}{\linewidth}
\centering
\normalsize
\setlength{\tabcolsep}{3pt}
\renewcommand{\arraystretch}{1.10}
\setlength{\belowcaptionskip}{\baselineskip}
\captionof{table}{Shared adaptation hyperparameters for SmartVLA-Pi and SmartVLA-OFT, listed by target domain and protocol. Seed denotes the training random seed.}\label{tab:domain-finetuning-hparams}
\begin{tabular}{@{}>{\raggedright\arraybackslash}p{0.300\linewidth}>{\raggedright\arraybackslash}p{0.170\linewidth}>{\raggedright\arraybackslash}p{0.170\linewidth}>{\raggedright\arraybackslash}p{0.160\linewidth}>{\raggedright\arraybackslash}p{0.130\linewidth}@{}}
\toprule
Domain / protocol & Updates & Global batch & Peak LR & Seed \\
\midrule
LIBERO & 50,000 & 256 & $10^{-5}$ & 42 \\
RoboTwin Full & 120,000 & 256 & $10^{-5}$ & 42 \\
RoboTwin C2R & 50,000 & 256 & $10^{-5}$ & 42 \\
Piper & 50,000 & 64 & $10^{-5}$ & 1000 \\
\bottomrule
\end{tabular}
\end{minipage}
\end{center}
%

\hypertarget{agentic-calling-protocol}{%
\subsection{Agentic Calling
Protocol}\label{agentic-calling-protocol}}

In Table~\ref{tab:libero-smart}, direct-policy results use all 6,000
instances (120 tasks, 10 macro layouts, and 5 local initializations per
layout). The agentic protocol uses the same task contracts, and model
comparisons are made within each protocol.

In an agentic system, one upper-level call specifies a natural-language
command and the number of action blocks to execute. A single low-level
VLA prediction is \(\mathbf{a}_{t:t+H-1} \sim \pi_\theta(\cdot \mid
\mathbf{o}_t,\mathbf{s}_t,\mathbf{c}_t)\),
where \(\mathbf{o}_t\) contains the current third-person and wrist RGB
observations. Colors, regions, trajectories, and target points involved
in the task appear directly in these images rather than through
privileged channels. \(\mathbf{s}_t\) denotes end-effector pose and
gripper state, and \(\mathbf{c}_t\) is the natural-language command
supplied by the caller. It may be a complete task, a current subgoal, or
a local corrective command based on execution. SmartVLA predicts a local
block of length \(H\); our Pi-based agentic evaluations use \(H=10\). This is the length of one
action
block, not the total number of environment steps in an upper-level call.

Each upper-level call also specifies the number \(K\) of action blocks,
corresponding to a budget of \(KH\) environment action steps. Within the
call, the VLA executes blocks under the same command: after each block,
the system obtains updated observations and robot state and predicts the
next block. After \(K\) blocks, control returns to the caller. Thus, one
upper-level call can contain multiple action blocks and still use
updated observations for closed-loop control; it is not required to
complete an entire command or long-horizon task.

The human operator or upper-level agent can use the returned observation
to call the VLA again with the original command, a sub-task command, or
a corrective command, and can choose a new block count \(K\). It may
instead call an atomic skill in the agentic system or terminate the
episode. Both \(K\) and the number of calls are selected dynamically
from execution progress and remain subject to a common interaction
budget.

For the Agentic+$\pi_{0.5}$ and Agentic+SmartVLA-Pi comparison, we use GPT-5.6
as the fixed upper-level planning model. We fix the upper-level
controller, available tools, calling protocol, and re-planning rules, and
keep tasks, initializations, and interaction budgets identical while
replacing only the low-level VLA. The upper controller decides from
observations when to call the VLA, which command to use, and how many
action blocks to execute, but cannot access environment rewards, online
success signals, or hidden evaluation state.

CaP-X uses its own harness and is reported as a separate complete-system
baseline rather than as a controlled replacement of the low-level VLA.

The agentic systems use the same scoring rules. An episode terminates when
the environment-step budget is exhausted or the upper-level VLM reports
task completion, whichever occurs first. A completion report is a stopping
decision, not an evaluator-provided success signal. The evaluator retains
the episode's contract history across VLA calls and tool invocations;
returning control to the upper level does not clear forbidden events or
the recorded step sequence. Success is determined by strict contract
satisfaction at termination, without the additional 30-step success window
used for end-to-end policies. A VLM completion report alone does not count
as success.

All agentic scores measure \textbf{overall Agent+VLA performance}, not the
standalone capability of the VLA. The fixed-controller comparison between
Agentic+$\pi_{0.5}$ and Agentic+SmartVLA-Pi isolates the effect of replacing the
low-level policy within that shared upper-level system.

